% Generated draft: requires mathematical and source review.
\documentclass[11pt]{article}
\usepackage[margin=25mm]{geometry}
\usepackage{fontspec}
\setmainfont{Latin Modern Roman}
\usepackage{amsmath,amssymb,mathtools,unicode-math}
\setmathfont{Latin Modern Math}
\usepackage{enumitem,needspace,xurl,hyperref}
\setlength{\parindent}{0pt}
\setlength{\parskip}{5pt}
\setlength{\emergencystretch}{3em}
\newcommand{\N}{\mathbb{N}}
\newcommand{\Z}{\mathbb{Z}}
\newcommand{\Q}{\mathbb{Q}}
\newcommand{\R}{\mathbb{R}}
\newcommand{\C}{\mathbb{C}}
\newcommand{\F}{\mathbb{F}}
\newcommand{\A}{\mathbb{A}}
\newcommand{\PP}{\mathbb{P}}
\newcommand{\E}{\mathbb{E}}
\newcommand{\eps}{\varepsilon}
\newcommand{\dd}{\,\mathrm{d}}
\begin{document}
\section*{UnsolvedMath ID: 30001334}
\textbf{Source problem code:} catalogue entry
\subsection{Definitions and mathematical statement}
For context, let \(A\) and \(B\) be \(C^*\)-algebras. In a standard convention, an \(A\)-\(B\) correspondence is a right Hilbert \(B\)-module \(E\), together with a \(*\)-homomorphism \(\varphi_E\colon A\to\mathcal L_B(E)\), where \(\mathcal L_B(E)\) denotes the algebra of adjointable \(B\)-linear operators on \(E\). Some conventions additionally require nondegeneracy, namely \(\overline{\operatorname{span}(\varphi_E(A)E)}=E\); the supplied record does not specify this requirement. An isomorphism between two such correspondences \(E\) and \(F\) is a surjective \(B\)-linear map \(U\colon E\to F\) satisfying
\[
\langle U\xi,U\eta\rangle_B=\langle\xi,\eta\rangle_B,
\qquad
U\varphi_E(a)\xi=\varphi_F(a)U\xi
\quad(a\in A,\ \xi,\eta\in E).
\]
In the usual correspondence framework, objects are \(C^*\)-algebras, \(1\)-morphisms are correspondences, and composition uses the interior tensor product: an \(A\)-\(B\) correspondence \(E\) and a \(B\)-\(C\) correspondence \(F\) compose to \(E\otimes_B F\). The supplied record does not specify the permitted \(2\)-morphisms. The ordinary tensor-product construction has associativity and unit identifications up to canonical isomorphism, so the term ``\(2\)-category'' may refer to a bicategory or a chosen strictification.

\textbf{Source question.} Why is it useful, in the correspondence \(2\)-category of \(C^*\)-algebras, to retain individual correspondences and isomorphisms between them instead of replacing correspondences by their isomorphism classes?

This is a conceptual request for an explanation. The supplied record gives no quantified conjecture, particular construction, universal property, or criterion of usefulness to establish. The definitions above provide a standard interpretive framework; they do not resolve the unspecified conventions or turn the question into a theorem.
\subsection{Short English statement}
The question asks why a correspondence \(2\)-category should retain isomorphic \(C^*\)-algebra correspondences as separate \(1\)-morphisms connected by isomorphisms. The intended benefit and precise categorical conventions are unspecified, so this remains a conceptual question requiring human review.
\subsection{Sources}
\begin{itemize}
\item Supplied local UnsolvedMath registry record: permanent ID \(30001334\), source problem code \texttt{catalogue entry}, title ``Isomorphism Data in Correspondence Two-Categories,'' in the Oberwolfach Reports Open Problems collection. No individual proposer is supplied; the proposed year is recorded as 2009.
\item Supplied source citation: \emph{Noncommutative Geometry (2010)}, Oberwolfach Reports, pp.\ 2271--2334, DOI \texttt{10.4171/owr/2009/41}. Original Oberwolfach report: \url{https://doi.org/10.4171/owr/2009/41}. The supplied citation names 2010, while the DOI and proposed-year metadata contain 2009; these details are retained without reconciliation. No problem-specific page or speaker attribution is supplied.
\item Supplied additional literature source: ``MFO report on C-star correspondences,'' \url{https://ems.press/content/serial-article-files/46243?nt=1}. The record supplies no further bibliographic details or pinpoint reference.
\item Literature assessment dated exactly \texttt{2026-08-21}: the supplied curation labels the question open, describes it as conceptual rather than a theorem with a determinate truth value, and states that later work explains the \(2\)-categorical role of retaining correspondence isomorphism classes. This is supplied dated curation, not independently verified. The phrase about retaining isomorphism classes requires clarification in light of the question's contrast with identifying isomorphic correspondences.
\item The supplied statement status is \texttt{corrected\_verified}; its verification note reports comparison with the official Oberwolfach report and typesetting of operator-algebra notation without changing the source-level claim. This is a supplied verification claim, not independent verification in this draft. The record expressly requests independent checking of time-sensitive status and cited claims.
\end{itemize}
\end{document}
